The standard \Sn equations approximate the continuous angular flux as a system of delta functions on the unit sphere.
These equations are then coupled together through the scattering source term, which is an integral over direction vector.
For a monoenergetic, time-independent problem, the system of equations has the form:
\begin{equation}
  \label{eq:sn-equations}
  \ohat_n\cdot\nabla\psirn+\Sigma_t\phaser\psirn=\qrn
\end{equation}
for $n\in N$ discrete directions (ordinates), which are preselected from a quadrature set on the unit sphere.
The source term, $\qrn$, represents all sources present, which can include:
\begin{enumerate}
  \item scattering
  \item fission
  \item multiplication not via fission ($(n,2n)$, $(n,3n)$)
  \item delayed neutron precursor decay
  \item independent volumetric source.
\end{enumerate}
For a non-multiplying medium, scattering and the independent source are the only source terms, and the latter does not change throughout a simulation.
The scattering source is defined as:
\begin{equation}
  \label{eq:scattering-source}
  q_{scatter}=\int_{4\pi}\Sigma_s\left(\rvec,\ohat'\cdot \ohat\right)\psi\left(\rvec,\ohat'\right)\odif{\ohat'}.
\end{equation}
In an \Sn formulation, the integral \cref{eq:scattering-source} is approximated via a quadrature set.
This set of discrete angles determines the angles for which \cref{eq:sn-equations} is solved for.

The benefits to the \Sn model include in simplistic formulation and good numerical properties, such as strong stability and good parallelization.
However, in replacing the standard model, a new numerical artifact presents itself: ray effects.

  [FILL IN MORE]

A new formulation is proposed, which we have called the \gls{wrldo}.
Beginning with the steady-state, monoenergetic \gls{bte}\footnote{In principle, the most general form of the \gls{bte} follows the same derivations.}:
\begin{equation}
  \label{eq:ss-me-bte}
  \ohat\cdot\nabla\psi\phasero+\Sigma_t\phaser\psi\phasero=q\phasero.
\end{equation}
Or, by expanding the streaming term:
\begin{equation}
  \label{eq:exp-ss-me-bte}
  \ohat_x\pdv{\psi}{x}+\ohat_y\pdv{\psi}{y}+\ohat_z\pdv{\psi}{z}+\Sigma_t\phaser\psi\phasero=q\phasero.
\end{equation}
We then perform two steps.
First, we approximate the angular flux with a spectral expansion in angle, where the basis functions are chosen to be Lagrange polynomials.
The direction vector, $\ohat$, can be uniquely defined by an azimuthal and polar angle.
For the set of $J+1$ azimuthal angles $\theta_j\in [\theta_0,\theta_1,\ldots,\theta_{J}]$ and the set of $I+1$ polar angles $\varphi_i\in [\varphi_0,\varphi_1,\ldots,\varphi_I]$, the Lagrange interpolants are built with a tensor product:
\begin{equation}
  \psi\phasero\approx\sum_{i=0}^I\sum_{j=0}^J \psi_{ij}\left(\rvec\right)\ell_i(\varphi) \ell_j(\theta).
\end{equation}
For clearer notation, we define a function $\mathcal{L}$ and a mapping to a flattened index $n=i(J+1)+j$:
\begin{equation}
  \mathcal{L}_n\left(\ohat\right) \equiv \ell_{\lfloor n/(J+1)\rfloor}(\varphi)\ell_{n \bmod (J+1)}(\theta)
\end{equation}
This lets us rewrite the spectral expansion as:
\begin{equation}
  \label{eq:spectral-exp-angle}
  \psi\phasero\approx\sum_{n=0}^N\psirn\mathcal{L}_n\left(\ohat\right)
\end{equation}
for $N+1 = (I+1)\times (J+1)$ unique direction vectors.
The nodes are chosen to correspond to the abscissae of a quadrature set on the unit sphere.
Second, we multiply \cref{eq:ss-me-bte} by a test function, $\gamma\left(\ohat\right)$, and integrate over the unit sphere.
This becomes:
\begin{multline}
  \sum_{n=0}^N\pdv{\psi_n}{x} \int_{4\pi}\ohat_x \mathcal{L}_n\left(\ohat\right)\gamma\left(\ohat\right) \odif{\ohat} +
  \sum_{n=0}^N\pdv{\psi_n}{y} \int_{4\pi}\ohat_y\mathcal{L}_n\left(\ohat\right)\gamma\left(\ohat\right)\odif{\ohat} + \\
  +\sum_{n=0}^N\pdv{\psi_n}{z} \int_{4\pi}\ohat_z\mathcal{L}_n\left(\ohat\right)\gamma\left(\ohat\right)\odif{\ohat} +
  \sum_{n=0}^N \Sigma_t \int_{4\pi}\psirn \mathcal{L}_n\left(\ohat\right)\gamma\left(\ohat\right)\odif{\ohat} = \\
  =\int_{4\pi}q\phasero\gamma\left(\ohat\right)\odif{\ohat}.
\end{multline}
Finally, we choose $\gamma\left(\ohat\right) \in [\mathcal{L}_0,\mathcal{L}_1,\ldots,\mathcal{L}_N]$, forming a Galerkin projection:
\begin{multline}
  \label{eq:wrldo-galerkin}
  \sum_{n=0}^N\pdv{\psi_n}{x} \int_{4\pi}\ohat_x \mathcal{L}_n\left(\ohat\right)\mathcal{L}_m\left(\ohat\right) \odif{\ohat} +
  \sum_{n=0}^N\pdv{\psi_n}{y} \int_{4\pi}\ohat_y\mathcal{L}_n\left(\ohat\right)\mathcal{L}_m\left(\ohat\right) \odif{\ohat} + \\
  +\sum_{n=0}^N\pdv{\psi_n}{z} \int_{4\pi}\ohat_z\mathcal{L}_n\left(\ohat\right)\mathcal{L}_m\left(\ohat\right) \odif{\ohat} +
  \sum_{n=0}^N \Sigma_t \int_{4\pi}\psirn \mathcal{L}_n\left(\ohat\right)\mathcal{L}_m\left(\ohat\right)\odif{\ohat} = \\
  =\int_{4\pi}q\phasero\mathcal{L}_m\left(\ohat\right) \odif{\ohat}.
\end{multline}
For clarity, we define the functional:
\begin{equation}
  \label{eq:sphere-functional}
  G_{mn}\left[f\left(\ohat\right)\right]=\int_{4\pi}f\left(\ohat\right)\mathcal{L}_n\left(\ohat\right)\mathcal{L}_m\left(\ohat\right) \odif{\ohat}
\end{equation}
and substitute this into \cref{eq:wrldo-galerkin}, giving:
\begin{multline}
  \sum_{n=0}^N \left[G_{mn}\left[\ohat_x\right]\pdv{\psi_n}{x}+G_{mn}\left[\ohat_y\right]\pdv{\psi_n}{y}+ G_{mn}\left[\ohat_z\right]\pdv{\psi_n}{z} + G_{mn}\left[1\right]\Sigma_t\psirn\right]=\\=\int_{4\pi}q\phasero\mathcal{L}_m\left(\ohat\right) \odif{\ohat}.
\end{multline}
We can now address the evaluation of the functional.
Approximating \cref{eq:sphere-functional} with a generic quadrature set gives:
\begin{equation}
  \begin{split}
    \int_{4\pi}f\left(\ohat\right)\mathcal{L}_n\left(\ohat\right)\mathcal{L}_m\left(\ohat\right) \odif{\ohat} & \approx \sum_{k=0}^N \rho_k f\left(\ohat_k\right) \underbrace{\mathcal{L}_n\left(\ohat_k\right)}_{\delta_{nk}} \underbrace{\mathcal{L}_m\left(\ohat_k\right)}_{\delta_{mk}} \\
                                                                                                              & =\rho_m f(\ohat_m)\delta_{mn}
  \end{split}
\end{equation}
where in the last line we have use the interpolation property of the Lagrange interpolants to filter out most terms in the sum.
This in term creates a diagonal system of equations $\underline{\Omega}\vec{\psi}=\vec{q}$, where:
\begin{equation}
  \label{eq:diagonal-system}
  \begin{aligned}
     & \underline{\Omega}_{mn} & = & \begin{cases}
                                       \rho_n\left[\left(\ohat_x\right)_m\pdv{}{x}+\left(\ohat_y\right)_m\pdv{}{y}+\left(\ohat_z\right)_m\pdv{}{z}+\Sigma_t\right], & m=n     \\
                                       0,                                                                                                                           & m\neq n
                                     \end{cases} \\
     & \vec{q}_m               & = & \int_{4\pi}q\phasero\mathcal{L}_m\left(\ohat\right) \odif{\ohat}.
  \end{aligned}
\end{equation}
Letting $m\leftarrow n$ let's us write the $n$\textsuperscript{th} equation of \cref{eq:diagonal-system} as:
\begin{equation}
  \label{eq:sn-like}
  \left(\ohat_x\right)_n\pdv{\psi_n}{x}+\left(\ohat_y\right)_n\pdv{\psi_n}{y}+\left(\ohat_z\right)_n\pdv{\psi_n}{z}+\Sigma_t\psi_n =  \mathcal{Q}_n
\end{equation}
where:
\begin{equation}
  \label{eq:lagrange-source}
  \mathcal{Q}_n\equiv\frac{1}{\rho_n}\int_{4\pi}q\phasero\mathcal{L}_n\left(\ohat\right)\odif{\ohat}.
\end{equation}
Those familiar with numerical transport methods may recognize \cref{eq:sn-like} as an \Sn-like equation with a Lagrange-weighted integral source term, termed the \textit{Lagrange source}.
\Cref{eq:sn-like} can thus be solved using swept methods, which are highly favorable for modern computer architecture.

Astonishingly, if one were to numerically compute the integral in \cref{eq:lagrange-source} using the \emph{same} quadrature set as the one over which the spectral expansion nodes of the angular flux (\cref{eq:spectral-exp-angle}) were selected, the Lagrange source collapses to the standard \Sn formulation.
However, the purpose of this derivation is to provide a new formulation that allows for a higher-order representation of the source term.
This will hopefully lead to improved accuracy and the reduction of numerical artifacts such as ray effects.
Therefore, going forward, we define a new, higher-order quadrature set on the unit sphere:
\begin{equation}
  \label{eq:high-order-spherical-quad-set}
  \int_{4\pi}f\left(\ohat\right)\odif{\ohat}\approx \sum_{p=0}^{P}w_p f\left(\ohat_p\right)
\end{equation}
for $P+1$ total quadrature points that will be used to evaluate integrals present in the source.

We now explore the evaluation of $\mathcal{Q}_n$.
For a multiplying system with an independent source, $q$ is a sum of the contributions,
\begin{equation}
  q\phasero = q_s\phasero + q_f\left(\rvec\right)+q_{ind}\phasero,
\end{equation}
where $q_s$, $q_f$, and $q_{ind}$ are the scattering, fission, and independent sources, respectively.
Note that the fission source is \emph{only} a function of space -- this is because fission is assumed to be an isotropic process.

Examining the scattering source and using \cref{eq:spectral-exp-angle}:
\begin{equation}
  \label{eq:scatter-source-no-quad}
  \begin{split}
    q_{s}\phasero & = \int_{4\pi}\Sigma_s\left(\rvec,\ohat'\cdot\ohat\right)\psi\left(\rvec,\ohat'\right)\odif{\ohat'}                                                                 \\
                  & \approx\int_{4\pi}\Sigma_s\left(\rvec,\ohat'\cdot\ohat\right)\left[\sum_{n=0}^N\psi_{n}\left(\vec{r}\right)\mathcal{L}_{n}\left(\ohat'\right)\right] \odif{\ohat'} \\
                  & =\sum_{n=0}^N\left[\int_{4\pi}\Sigma_s\left(\rvec,\ohat'\cdot\ohat\right)\mathcal{L}_{n}\left(\ohat'\right)\odif{\ohat'}\right]\psi_n\left(\rvec\right).
  \end{split}
\end{equation}
Then, by applying \cref{eq:high-order-spherical-quad-set}, the integral above becomes:
\begin{equation}
  \int_{4\pi}\Sigma_s\left(\rvec,\ohat'\cdot\ohat\right)\mathcal{L}_{n}\left(\ohat'\right)\odif{\ohat'}  \approx \sum_{p=0}^P w_p \Sigma_s\left(\rvec,\ohat_p\cdot \ohat\right)\mathcal{L}_{n}\left(\ohat_p\right).
\end{equation}
We can go one step further and expand the differential scattering cross section using Legendre polynomials~\cite{larsenAdvancesDiscreteOrdinatesMethodology2010}.
Defining $\mu_p\equiv \ohat_p\cdot\ohat$, this is~\cite{bellNuclearReactorTheory1970}:
\begin{equation}
  \label{eq:scatter-xs-expansion}
  \begin{aligned}
     & \Sigma_s\left(\rvec,\mu_p\right) & = & \sum_{l=0}^{\infty}\Sigma_{s,l}\left(\rvec\right)P_l\left(\mu_p\right)                       \\
     & \Sigma_{s,l}                     & = & \frac{2l+1}{2}\int_{-1}^{1}\Sigma_s\left(\rvec,\mu_p\right)P_l\left(\mu_p\right)\odif{\mu_p}
  \end{aligned}
\end{equation}
Thus, \cref{eq:scatter-source-no-quad} is now:
\begin{equation}
  \label{eq:wrldo-scatter-source}
  q_s\phasero\approx\sum_{n=0}^N\psi_{n}\left(\rvec\right)\sum_{p=0}^{P}w_p\mathcal{L}_{n}\left(\ohat_p\right)\sum_{l=0}^{\infty}\Sigma_{s,l}P_l\left(\mu_p\right).
\end{equation}
In practice, the expansion in \cref{eq:scatter-xs-expansion} is truncated at a finite order, $L$, and the moments are computed before the simulation is run.

The fission source is much simpler to address:
\begin{equation}
  \label{eq:wrldo-fission-source}
  \begin{aligned}
    q_f & =\frac{1}{4\pi}\nu\Sigma_f\int_{4\pi}\psi\phasero\odif{\ohat}                                                                                                                                      \\
        & \approx \frac{1}{4\pi}\nu\Sigma_f \int_{4\pi}\left[\sum_{n=0}^N\psi_n\left(\rvec\right)\mathcal{L}_n\left(\ohat\right)\right]\odif{\ohat} & \text{(insert \cref{eq:spectral-exp-angle})}           \\
        & =\frac{1}{4\pi}\nu\Sigma_f \sum_{n=0}^N\psi_n\left(\rvec\right)\int_{4\pi}\mathcal{L}_n\left(\ohat\right)\odif{\ohat}                                                                              \\
        & \approx \frac{1}{4\pi}\nu\Sigma_f \sum_{n=0}^N \psi_n\left(\rvec\right)\left[\sum_{p=0}^P w_p\mathcal{L}_n\left(\ohat_p\right) \right]    & \text{(apply \cref{eq:high-order-spherical-quad-set})} \\
  \end{aligned}
\end{equation}
Notice in that applying our approximations, we have numerically defined the scalar flux:
\begin{equation}
  \label{eq:numerical-scalar-flux}
  \begin{split}
    \phi(\rvec) & \equiv\int_{4\pi}\psi\phasero\odif{\ohat}                                                         \\
                & \approx \sum_{n=0}^{N}\psi_n\left(\rvec\right)\sum_{p=0}^{P}w_p\mathcal{L}_n\left(\ohat_p\right).
  \end{split}
\end{equation}

Finally, we apply \cref{eq:high-order-spherical-quad-set} directly to \cref{eq:lagrange-source} to get a final expression for the Lagrange source term:
\begin{equation}
  \label{eq:discrete-lagrange-source}
  \mathcal{Q}_n \approx \frac{1}{\rho_n}\sum_{p=0}^{P}w_p q\left(\rvec,\ohat_p\right)\mathcal{L}_n\left(\ohat_p\right),
\end{equation}
where the source at each high-order quadrature node, $q\left(\rvec,\ohat_p\right)$, is evaluated using \cref{eq:wrldo-scatter-source}, \cref{eq:wrldo-fission-source}, and the independent source.

A source iteration algorithm is proposed and is illustrated in \cref{fig:wrldo-flowchart}.
The algorithm solves the iteration scheme:
\begin{subequations}
  \label{eq:wrldo-source-iters}
  \begin{equation}
    \left(\ohat_x\right)_n\pdv{\psi^{k+1}_n}{x}+\left(\ohat_y\right)_n\pdv{\psi^{k+1}_n}{y}+\left(\ohat_z\right)_n\pdv{\psi^{k+1}_n}{z}+\Sigma_t\psi^{k+1}_n  =  \mathcal{Q}^k_n
  \end{equation}
  \begin{equation}
    \phi^{k+1}\left(\rvec\right)=\sum_{n=0}^{N}\psi^{k+1}_n\left(\rvec\right)\sum_{p=0}^{P}w_p\mathcal{L}_n\left(\ohat_p\right)
  \end{equation}
  \begin{multline}
    q^{k+1}\left(\rvec,\ohat_p\right) = \sum_{n=0}^N\psi^{k+1}_{n}\left(\rvec\right)\sum_{p'=0}^{P}w_{p'}\mathcal{L}_{n}\left(\ohat_{p'}\right)\sum_{l=0}^{L}\Sigma_{s,l}P_l\left(\ohat_{p'}\cdot\ohat_p\right) + \frac{1}{4\pi}\nu\Sigma_f \phi^{k+1}\left(\rvec\right) +\\+q_{ind}\left(\rvec,\ohat_p\right)
  \end{multline}
  \begin{equation}
    \mathcal{Q}^{k+1}_n =  \frac{1}{\rho_n}\sum_{p=0}^{P}w_p q^{k+1}\left(\rvec,\ohat_p\right)\mathcal{L}_n\left(\ohat_p\right)
  \end{equation}
\end{subequations}

A few notes about the proposed procedure:
\begin{enumerate}
  \item The spectral expansion in \cref{eq:spectral-exp-angle} is expected to be ``good enough'' so that the solution between points on the initial, lower-order spherical quadrature set is well-approximated.
  \item The additional quadrature set, \cref{eq:high-order-spherical-quad-set}, is intended to be a higher order to better approximate the integrals in the Lagrange source, scattering source, and scalar flux.
  \item The order of the Legendre expansion of the scattering cross section, $L$, will need to be chosen ahead of time -- an infinite sum is not feasible.
  \item Any discretization in space is fair game -- let's do the \gls{sem}!
  \item The above derivation is easily extended to time- and energy-dependent simulations. No approximations have been made on these domains.
\end{enumerate}

Some programming notes:
\begin{enumerate}
  \item Scattering moments must be computed ahead of time.
  \item Evaluation of the Lagrange interpolants at the higher-order quadrature set nodes need only be done once, then their values stored for later reference.
\end{enumerate}

\begin{figure}
  \centering
  \begin{tikzpicture}[node distance=2.2cm]
  \node (start) [startstop] {Start};
  \node (input) [io, below of=start] {User simulation input.};
  \node (init_quad) [process, below of=input, xshift=-1.5cm] {Create spherical quadrature sets.};
  \node (initialize) [process, right of=init_quad, xshift=1.5cm] {Initialize angular flux and Lagrange source.};
  \node (sweep) [process, below of=input, yshift=-3cm] {Sweep solution across spatial domain.};
  \node (source_moments) [process, below of=sweep, yshift=-1cm] {Compute new Lagrange source.};
  \node (moments_converged) [decision, right of=source_moments, xshift=3cm] {Lagrange source has converged?};
  \node (k_update) [process, below of=moments_converged, yshift=-1.75cm] {Compute new \keff.};
  \node (k_converged) [decision, below of=k_update, yshift=-1cm] {\keff has converged?};
  \node (recompute_moments) [process, right of=k_converged, xshift=3cm] {Recompute Lagrange source.};
  \node (stop) [startstop, below of=k_converged, yshift=-1cm] {Stop};

  \draw [arrow] (start) -- (input);
  \draw [arrow] (input) -- (init_quad);
  \draw [arrow] (input) -- (initialize);
  \draw [arrow] (init_quad) -- (sweep);
  \draw [arrow] (initialize) -- (sweep);
  \draw [arrow] (sweep) -- (source_moments);
  \draw [arrow] (source_moments) -- (moments_converged);
  \draw [arrow] (moments_converged) -- node[anchor=west] {yes} (k_update);
  \draw [arrow] (k_update) -- (k_converged);
  \draw [arrow] (k_converged) -- node[anchor=north] {no} (recompute_moments);
  \draw [arrow] (recompute_moments) |- (sweep);
  \draw [arrow] (k_converged) -- node[anchor=west] {yes} (stop);
  \draw [arrow] (moments_converged) |- node[anchor=east,yshift=-0.5cm] {no} (sweep);
\end{tikzpicture}
  \caption{\Gls{wrldo} numerical solution flowchart.}
  \label{fig:wrldo-flowchart}
\end{figure}